\documentclass[../../proceedings/main.tex]{subfiles}
\begin{document}
% Bearbeite den Inhalt hier; alle Metadaten stehen in metadata.tex.
% Das Layout ist zentral in proceedings/ definiert.
% Orientierung fuer insgesamt hoechstens 4 Seiten, inklusive References:
% Abstract 5--8 Zeilen; Introduction 0.5 S.; State of the Art 0.5 S.;
% Method 1--1.25 S.; Results & Evaluation 0.75--1 S.;
% Discussion & Limitations 0.5 S.; Conclusion 0.25 S.
% Dies sind Richtwerte, keine erzwungenen Seitenumbrueche.
\begin{refsection}[\subfix{references.bib}]
\input{\subfix{metadata.tex}}
% Titel, Autoren, Abstract, Keywords und Projekt-URL kommen aus metadata.tex.
\PaperHeading

\section{Introduction \& Motivation}
% Ca. 0.5 Seite: Welches konkrete Problem bearbeitet ihr, warum ist es relevant,
% wer profitiert, und was ist euer Ziel bzw. eigener Beitrag? Kommt schnell zum
% eigentlichen Problem statt mit einer allgemeinen KI-Einfuehrung zu beginnen.
State the project question and the expected benefit for a specific audience.
End with a precise description of your own contribution.

\section{State of the Art}
% Ca. 0.5 Seite: Ordnet nur fuer euer Problem relevante Arbeiten, Systeme und
% Alternativen ein. Vergleicht Vor- und Nachteile und zeigt die Luecke auf,
% die eure Methode adressiert. Keine allgemeine Geschichte der KI.
Compare a relevant prior approach with an alternative and explain why neither
fully solves your problem. Cite the actual source, for example
\autocite{vaswani2017attention}; replace this illustrative citation if it does
not fit your project.

\section{Method}
% Ca. 1--1.25 Seiten: Beschreibt Architektur, Daten, Modelle, Algorithmen,
% Verarbeitung und Implementierung so, dass der Ansatz nachvollziehbar ist.
% Begruendet Entscheidungen: Warum dieses Modell/diese Datenquelle/dieser
% Ablauf? Welche Alternative und welche Trade-offs gab es? Keine Tool-Liste.
Describe the data flow shown in Figure~\ref{fig:team-workflow} and explain
why each step is needed. Replace the placeholder with your own local figure.
\begin{figure}[t]
  \centering
  % Fuer eine echte Grafik: \includegraphics[width=\columnwidth]{\subfix{figures/my-figure.pdf}}
  \fbox{\parbox{0.8\columnwidth}{\centering Input $\rightarrow$ Model $\rightarrow$ Output}}
  \caption{Illustrative processing pipeline; replace with your system.}
  \label{fig:team-workflow}
\end{figure}
Explain how inputs become outputs and where key design choices apply.

\section{Results \& Evaluation}
% Ca. 0.75--1 Seite: Zeigt, was tatsaechlich erreicht wurde und auf welcher
% Grundlage ihr die Qualitaet beurteilt (Testdaten, Baseline, Metrik, Setting).
% Nutzt messbare Werte oder systematische qualitative Befunde; "worked well"
% ist keine Evidenz. Kennzeichnet illustrative Zahlen und echte Messungen klar.
Define your metric before reporting it. Equation~\ref{eq:team-precision}
illustrates precision; choose a metric that fits your own task.
\begin{equation}
  \mathrm{precision}=\frac{\text{relevant retrieved items}}{\text{retrieved items}}.
  \label{eq:team-precision}
\end{equation}
Table~\ref{tab:team-results} shows how to compare a baseline with your method.
Replace the illustrative values with evidence from your evaluation.
\begin{table}[t]
  \centering
  \caption{Illustrative comparison; replace with measured results.}
  \label{tab:team-results}
  \begin{tabular}{lc}\toprule
    Method & Precision \\ \midrule
    Baseline & 0.60 \\
    Proposed & 0.75 \\ \bottomrule
  \end{tabular}
\end{table}

\section{Discussion \& Limitations}
% Ca. 0.5 Seite: Interpretiert Erfolge und Fehler kritisch. Unter welchen
% Bedingungen versagt der Ansatz? Welche Daten-, Bias-, Fairness-, Datenschutz-
% oder Sicherheitsprobleme bestehen? Wie weit lassen sich Ergebnisse uebertragen?
% Nennt konkrete Verbesserungen statt nur pauschal "more research is needed".
Discuss at least one failure case and one threat to the validity of your
evaluation. Explain the practical consequences of those limits.

\section{Conclusion}
% Ca. 0.25 Seite: Fasst Problem, Umsetzung, wichtigste Erkenntnis und einen
% sinnvollen naechsten Schritt sehr kurz zusammen. Keine neuen Ergebnisse.
Summarize the contribution and its main evidence in a few sentences, then name
one realistic next step.

% References zaehlen zum Vier-Seiten-Limit. Jeder zitierte Eintrag gehoert in
% die eigene references.bib; verwendet eindeutige Label-Praefixe pro Team.
\printbibliography[heading=subbibliography,title={References}]
\end{refsection}
\end{document}
